% answers_short.tex —— 2011 年·湖南卷·理科 · 简短答案速查
% 由 tools/make_short_answers.py 自 answers.tex 的 \gksolhead/\gkanswer 提取，
% 并与 2011/2011_数学真题_湖南理_答案要点.md 逐题比对后生成，勿手改（改源 MD 后重跑）。

\begin{gknotice}
\textbf{使用说明}：本节只给最终答案便于快速核对，完整推导见解析版。答案为编写者独立推导，\textbf{非官方标准答案}。
\end{gknotice}

\gkpart{选择题}
\begin{center}\begin{tabular}{ccccccccc}
\toprule
题号 & 1 & 2 & 3 & 4 & 5 & 6 & 7 & 8 \\
\midrule
答案 & D & A & A & C & C & D & A & D \\
\bottomrule
\end{tabular}\end{center}

\gkpart{填空题}
\begin{itemize}[leftmargin=2.6em]
  \item[\textbf{9.}] $2$。
  \item[\textbf{10.}] $9$。
  \item[\textbf{11.}] $\dfrac{2\sqrt3}3$。
  \item[\textbf{12.}] $25$。
  \item[\textbf{13.}] $\dfrac23$。
  \item[\textbf{14.}] $-\dfrac14$。
  \item[\textbf{15.}] $\dfrac2\pi$；$\dfrac14$。
  \item[\textbf{16.}] (1) $2$；(2) $1093$
\end{itemize}

\gkpart{解答题}
\begin{itemize}[leftmargin=2.6em]
  \item[\textbf{17.}] (1) $C=\dfrac\pi4$；(2) 最大值 $2$，此时 $A=\dfrac\pi3$，$B=\dfrac{5\pi}{12}$
  \item[\textbf{18.}] (1) $\dfrac3{10}$；(2) $P(X=2)=\dfrac14$，$P(X=3)=\dfrac34$，$E(X)=\dfrac{11}4$
  \item[\textbf{19.}] (1) 证明见解析；(2) $\dfrac{\sqrt{10}}5$
  \item[\textbf{20.}] (1) $y=\dfrac{5(3|v-c|+10)}v$；(2) $0<c\leqslant\dfrac{10}3$ 时 $v=10$；$\dfrac{10}3<c\leqslant5$ 时 $v=c$
  \item[\textbf{21.}] (1) $C_1:\dfrac{x^2}4+y^2=1$，$C_2:y=x^2-1$；(2) ① 证明见解析；② 存在，$k=\pm\dfrac32$
  \item[\textbf{22.}] (1) $2$ 个零点；(2) 证明见解析，可取 $M=\max\{a,2\}$
\end{itemize}
